\documentclass[12pt,a4paper,openany]{book}
% Preamble for AMACSS Constitution
% Package imports
\usepackage[utf8]{inputenc}
\usepackage[margin=1in]{geometry}
\usepackage{enumitem}
\usepackage{hyperref}
\usepackage{titlesec}

% Section formatting
\titleformat{\section}{\bf\Large}{Article \Roman{section}: }{0em}{}

% List formatting
\setlist[enumerate,1]{label=\thesection.\arabic*, ref=\thesection.\arabic*}
\setlist[enumerate,2]{label=\theenumi.\arabic*, ref=\theenumi.\arabic*}
\setlist[enumerate,3]{label=\theenumii.\arabic*, ref=\theenumii.\arabic*}
\setlist[enumerate,4]{label=\theenumiii.\arabic*, ref=\theenumiii.\arabic*}
\setlist[enumerate,5]{label=\theenumiv.\arabic*, ref=\theenumiv.\arabic*}


\begin{document}
\title{Constitution of The Association of Mathematical and Computer 
Science Students}
\date{Updated as of August 20th, 2025}
\maketitle

\vspace{1cm}

\chapter{Name and Purpose}

\section{Organization Name}

\subsection{Name}
The official name of the Campus Group shall be The Association of 
Mathematical and Computer Science Students. The Campus Group 
shall be referred to as the \quotes{Association}.

\subsection{Acronym}
The Association may be referred to by the acronym AMACSS.


\section{Organization Purpose}

\subsection{Non-profit Status}
The Association is a non-profit, student-run group that operates 
independently of any faculty, department, or administrative unit of 
the University of Toronto, to be referred as the \quotes{University}.

\subsection{University Policy Compliance}
The Association shall operate in accordance to the 
\emph{Policy on the Recognition of Student Groups}\textemdash 
to be refered to as the \quotes{University Policy}\textemdash 
or any statute, which may be substituted therefore, as amended from
time to time by the University of Toronto Governing Council.

\subsection{Departmental Student Association}
The Association shall be the official departmental student
association of the Department of Computer and Mathematical Sciences 
at the University of Toronto Scarborough\textemdash 
to be referred to as the \quotes{CMS department}\textemdash 
as a registered student group of the University.

\subsection{Mission}
The purpose, objectives, and mission of the Association is to 
represent, advocate for and enhance the university experience of 
students in the CMS department\textemdash to be referred to as
\quotes{CMS students}\textemdash through the following:

\begin{enumerate}

    \item To offer academic and professional support and provide 
    resources to CMS students,

    \item To foster a sense of community among CMS students 
    through various social initiatives and activities,

    \item To promote collaboration and networking opportunities among 
    University students, recognized student groups, faculty, staff, 
    and industry, and

    \item To advocate for the interests and concerns of CMS students 
    to the university administration and relevant bodies.

\end{enumerate}


\chapter{Membership}

\section{Types of Membership}

\subsection{Memberships Offers}
The Association shall offer three (3) types of memberships that are not 
necessarily mutually exclusive:
\begin{enumerate}

    \item General (voting) membership,

    \item Honorary (non-voting) membership, and

    \item Paid (non-voting) membership.

\end{enumerate}

\subsection{Terms of Membership}
The term of all memberships shall start on May 1st and end on April 
30th of the following year.

\subsection{Membership List}
The Association shall maintain a list of members of the categories
and ensure that membership is being offered in accordance with
this constitution for every academic year.

\section{General Membership}

\subsection{General Membership Eligibility}
The eligible for general (voting) membership shall be as all UTSC
students.

\subsection{General Membership Termination}
General (voting) membership of the Association shall be 
automatically terminated if a member ceases to be a registered
student at the University of Toronto Scarborough
 or withdraws their membership.

\subsection{General Membership Fee}
There shall be no fee for general (voting) membership.


\subsection{General Membership Rights}
General members shall have the rights to:
\begin{enumerate}

    \item Have one vote in Association referenda, and
    Members Meetings, and

    \item Sign petitions related to the Association

\end{enumerate}

\section{Honorary Membership}

\subsection{Honorary Membership Purpose}
The purpose of honorary membership is to recognize
and honour individuals who have demonstrated exceptional dedication
and service to the Association or the CMS department. The membership
is intended to acknowledge faculty, staff, alumni, or community members
who have positively impacted the Association or the CMS students.

\subsection{Honorary Membership Eligibility}
Honorary membership may be granted to any individual 
who has made significant contributions to the Association or the 
CMS department, or be nominated and approved by the Board of 
Executives.

\subsection{Honorary Membership Termination}
Honorary membership may be revoked by a two-thirds
(2/3) majority vote of the Board of Executives. Honorary members
may also voluntarily withdraw their membership at any time.

\section{Paid Membership}

\subsection{Paid Membership Purpose}
The purpose of paid membership is to provide additional
benefits and services to members who choose to support the Association
financially. Paid members may receive exclusive access to certain events,
resources, or discounts on merchandise, as determined by the Board of
Executives.

\subsection{Paid Membership Eligibility}
Paid membership may be offered to individuals who may
already be general members or honorary members,
or to individuals who do not meet the criteria for general or honorary
membership.

\subsection{Paid Membership Fee}
The fee for paid membership shall be determined by the
Board of Executives and may be subject to change on an annual basis.

\subsection{Paid Membership Termination}
Paid membership shall be automatically terminated if a
member fails to pay the membership fee within the specified time frame
or withdraws their membership.

\subsection{Paid Membership Rights}
Paid members shall have the rights to access the benefits
and services associated with their membership, as determined by the
Board of Executives for the academic year.
However, they shall not have voting rights unless they are also
a general member.


\chapter{Board of Executives}

\section{Board of Executives Composition}

\subsection{Board Purpose}
The Board of Executives\textemdash to be referred to as the 
\quotes{Board}\textemdash 
shall be the governing body of the Association and shall be responsible 
for the overall management and direction of the Association. 
The Board shall represent the interests of the members and serve the 
Association in accordance with this constitution, policies, and its
constituents.

\subsection{Term of Office}
The term for all positions on the Board shall be from 
May 1st to April 30th.

\subsection{Composition}
The Board shall be comprised of two (2) comittes: 
Senior Executives Committee and Executive Committee. 
These committee positions may hold voting or non-voting positions within 
the Board as specified.

\section{Board of Executives Eligibility and Restrictions}

\subsection{Voting Executives Eligibility}\label{voting-exec-eligibility}
All voting members of the Board must be currently 
registered students of the University of Toronto.

\subsection{Non-voting Executives Eligibility}\label{non-voting-exec-eligibility}
Non-voting members may hold only non-voting positions on the Board.
These positions may be held by individuals who are not
currently registered students of the University of Toronto, such as
faculty, staff, alumni, or community members.

\subsection{Not Registered Student Executive Positions}
The maximum amount of non-voting positions held by those who are not
registered students on the Board shall be one (1) position or ten 
per cent (10\%) of the positions on the Board, whichever is greatest.

\subsection{Restrictions on Non-voting Executive}
Persons holding non-voting positions on the Board cannot serve as an 
officer, financial authority, signing authority, primary contact, or 
secondary contact.

\subsection{One Voting Executive Incumbency Policy}
There shall be no individual who may hold more than one (1) voting 
position on the Board at any given time unless they are serving as
acting capacity for an unfilled position, to which they still only
hold one (1) vote in the Board. There also shall only 
be one (1) incumbent per voting position on the Board at any given 
time.
 
\subsection{Financial Authority Restrictions}
No person may serve as a financial authority or signing authority 
for the Association if they are currently serving as a financial 
authority or signing authority for another recognized student group 
at the University of Toronto.

\subsection{Appointment of Directors or Coordinators}
The Board of Executives may appoint Directors or Coordinators for 
various committees who do not hold executive decision-making 
authority and are not eligible to cast votes at meetings of the 
Board.

\section{Senior Executive Committee}

\subsection{Senior Executive Positions}
The Senior Executive Committee shall consist of the following four
(4) voting positions elected by the Association general members:
\begin{enumerate}

    \item President,
    \item Senior Vice-President Communications,
    \item Senior Vice-President Academic Affairs, and
    \item Senior Vice-President Operations.
    

\end{enumerate}

\subsection{Senior Executive Purpose}
The Senior Executive Committee shall be the chief governing body
of the Association and shall be responsible for the overall strategic
management and direction of the Association. Each Senior Vice-President
shall oversee a specific division of the Association and are directly
responsible for the strategic operations and future growth of their 
respective divisions.

\subsection{Non-voting Senior Executive Positions}
The Senior Executive Committee shall additionally have one (1)
Chief of Staff as a non-voting position appointed by the Senior
Executive Committee.

\subsection{Senior Executive Exclusion Policy}
No Senior Executive Committee member may hold any other executive
position from any other recognized student group at the University of
Toronto during their incumbency.

\subsection{President} The President shall be the chief executive officer 
over internal organization and the figurehead representing all members
of the Association. The President shall:
\begin{enumerate}

    \item Oversee and manage the organization, maintain the 
    integrity of the Association, and ensure all events and functions
    align with the mission of the Association,

    \item Act as the primary representative and spokesperson for the
    Association in all official matters, including interactions with
    University offices, the Scarborough Campus Student Union,
    other recognized student groups, and external organizations,
    and the public as a Departmental Student Association,

    \item Lead the Senior Executive Committee in strategic planning
    and decision-making processes of the Association accross all
    sectors,

    \item Conduct meetings of the Board of Executives and general 
    members, ensuring effective communication and collaboration among 
    all members of the team,

    \item Write executive orders and ensure the enforcement of the
    constitution, policies, and procedures of the Association, and,

    \item Hold special veto privileges over decisions made by the Board of
    Executives as specified.

\end{enumerate}

\subsection{Senior Vice-President Communications}
The Senior Vice-President Communications shall be the chief
communications officer of the Association and oversee strategic
development and operations of the Communications Division. 
The Senior Vice-President Communications shall:
\begin{enumerate}

    \item Develop and implement communication strategies to effectively
    reach and engage CMS students, promote Association events, disseminate
    important information to University bodies and external organizations,

    \item Manage the Association's public precense with general members,
    University bodies, recognized student groups, and external organizations,

    \item Act as the first point of contact for the
    Association in internal department affairs, such as general members,
    CMS recognized student groups, and the CMS department, and 

    \item Direct and supervise the vice-presidents and other team 
    members of the Communications Division.

\end{enumerate}

\subsection{Senior Vice-President Academic Affairs}
The Senior Vice-President Academic Affairs shall be the chief
academics officer of the Association and oversee stragic development
and operations of the Academic Division. 
    The Senior Vice-President Academic Affairs shall:
\begin{enumerate}

    \item Liaise with faculty members, general members, academic 
    advisors, and other relevant stakeholders to address academic 
    concerns and advocate for the needs of CMS students,

    \item Develop, implement, and maintain academic support programs,
    workshops, and resources to enhance the academic experience of 
    CMS students according to the mission of the Association, and

    \item Direct and supervise the vice-presidents and other team 
    members of the Academic Division to serve the academic needs of 
    the CMS students.

\end{enumerate}

\subsection{Senior Vice-President Operations}
The Senior Vice-President Operations shall be the chief operations
officer of the Association and oversee strategic development and
operations of the Operations Division. 
The Senior Vice-President Operations shall:
\begin{enumerate}

    \item Oversee the planning, coordination, and execution of
    events, initiatives, and services provided by the Association
    to CMS students,

    \item Manage logistical aspects of the Association's activities,
    including venue bookings, resource allocation, and volunteer
    coordination, and

    \item Direct and supervise the vice-presidents and other team 
    members of the Operations Division.

\end{enumerate}

\subsection{Chief of Staff}
The Chief of Staff shall be the primary aide to the Senior Executive
Committee and assist in the coordination and management of the Board
of Executives and the Association internal conduct.
The Chief of Staff shall:
\begin{enumerate}

    \item Assist the Senior Executive Committee in strategic planning,
    decision-making, and execution of the Association's mission,

    \item Ensure the Board of Executives operates efficiently, effectively,
    and in accordance with the Association's constitution, policies and 
    procedures,

    \item Advise executive members on best practices, governance, and
    organizational improvements,

    \item Moderate meetings of the Board of Executives and chair member
    meetings, and

    \item Manage special projects, initiatives, and tasks as assigned
    by the Senior Executive Committee.

\end{enumerate}

\section{Executive Committee}

\subsection{Executive Committee Positions}
The Executive Committee offices shall consist of the following
(10) voting positions appointed by the Senior Executive
Committee:
\begin{enumerate}

    \item Vice-President Academics
    \item Vice-President Review Seminars
    \item Vice-President Finance
    \item Vice-President Logistics
    \item Vice-President Lounge
    \item Vice-President Special Events
    \item Vice-President University Affairs
    \item Vice-President Sponsorship
    \item Vice-President Engagement
    \item Vice-President Marketing

\end{enumerate}

\subsection{Interim Vice-Presidents}
There may be interim vice-presidents
appointed by the Senior Executive Committee for the following:
\begin{enumerate}
    
    \item A vice-president position is vacant, or
    
    \item The Senior Executive Committee deems it necessary to appoint
    a person who does not meet the eligibility requirements as stated
    in §\ref{voting-exec-eligibility} for their expertise and
    value they can contribute to the association. This interim
    vice-president shall be a non-voting member on the Board to
    comply with the University Policy.
\end{enumerate}


\subsection{Coordinators}
The Executive Committee may additionally have non-voting executive
positions called Coordinators that are appointed by the Senior 
Executive Committee to assist with various teams.

\subsection{Faculty Advisor}
The Executive Committee shall additionally have one (1) Faculty
Advisor as a non-voting position appointed by the Senior Executive
Committee in consultation with the CMS department.


\subsection{Advisory Group}
The Executive Committee shall additionally have an Advisory Group
containing honorary members appointed by the Senior Executive 
Committee to provide guidance and support to the Board.
The Advisory Group may include alumni, faculty, staff, or community
members who have demonstrated a commitment to the Association and
its mission.

\subsection{Executive Advisor}
The Executive Committee shall additionally have non-voting executive


\chapter{Elections of Senior Executive}

\section{Election Logistics}

\subsection{Election Timing}
Elections for all voting positions on the Senior Executive Committee 
shall be held annually in the month of February and to be completed by the 
end of March.

\subsection{Election Process}
All voting positions on the Executive Committee shall be filled 
through an annual election. All voting-members of the Board
shall have one (1) vote for each position on the Senior Executive
Committee.

\subsection{Candidate Eligibility}
The eligible for all candidates on the Senior Executive 
Committee\textemdash in addition to the requirements of 
§\ref{voting-exec-eligibility}\textemdash shall be:
\begin{enumerate}

    \item A general members of the Association,
    
    \item A current member of the Board in the current term, or
    
    \item Nominated by one (1) member of the Board, or
    
    \item Nominated by twenty five (25) general members of the Association.

\end{enumerate}

\subsection{Election Winners}
The nominee winning the simple majority of votes cast in the election for 
each position shall be deemed the winner. In an event that no candidate
receives a simple majority of votes cast, a runoff election shall be held
between the two (2) candidates with the highest number of votes.

\end{document}
